% An echo-type sequence applied to the root pulse partially mitigates decoherence effects but remains limited by pronounced amplitude-dependent broadening. Conversely, the Lorentzian echo variant demonstrates an amplitude-independent linewidth spanning over three orders of magnitude in drive strength, evidencing remarkable robustness against amplitude-induced dephasing. This insensitivity arises from the adiabatic evolution and the symmetric frequency-splitting intrinsic to the pulse design. The primary limitation lies in the floating offset and amplitude-dependent baseline, resulting in reduced contrast compared to conventional square and Lorentzian drives.



To identify the useful operating regime, we examine the linewidth and contrast
of the echo response as functions of pulse duration and drive amplitude.  We
choose pulse durations shorter than $T_1$, but long enough that Fourier
broadening does not dominate.  The duration-dependent linewidth, contrast,
and contrast-weighted resolution are reported in the Supplemental
Material~\cite{SupplementalMaterial}.

A second design constraint is the cutoff of the formally infinite-support
pulse envelope.  Truncation introduces a residual endpoint drive
$\Omega_c=c\Omega_0$ and a corresponding cutoff-broadening
contribution~\cite{boradjiev_qubit_2013,mihov_defying_2024,mihov_qubit_2024}.
The relevant weak-saturation condition is
$s_c=\Omega_c^2T_1T_2\ll1$, or
$\Omega_c\ll1/\sqrt{T_1T_2}$; the inequality therefore points toward
\emph{smaller}, not larger, endpoint amplitudes.  The relative cutoff $c$
cannot be specified independently of $\Omega_0$, $T_1$, and $T_2$, giving a
three-dimensional design space of pulse duration, drive amplitude, and cutoff
level.  The corresponding cutoff--amplitude measurements and their
resolution--signal analysis are reported in the Supplemental
Material~\cite{SupplementalMaterial}.

Beyond linewidth and contrast, amplitude robustness requires that increasing
the drive not displace the inferred resonance center.  At large amplitudes,
multilevel AC-Stark shifts and finite-pulse interference become important; we
therefore examine the center stability in a measurement on a second qubit
over a calibrated nominal Rabi-frequency range of $10$--$60~\mathrm{MHz}$.

\begin{figure}[t]
    \centering
    \includegraphics[width=\columnwidth]{%
        ../figures/paper/04_main_drag_summary.pdf}
    \caption{DRAG-assisted echo-root-Lorentzian spectroscopy.
    Measured populations for (a) $\beta=0$ and (b) $\beta=-0.22$,
    with $L=50~\mu\mathrm{s}$ and cutoff $c=0.00075$; white curves
    trace fitted centers after subtracting each run's mean, lightly
    smoothed with a Gaussian of one amplitude-sample standard deviation.
    (c) The same smoothed center traces; the unsmoothed fits have RMS
    values of $6.78$ and
    $2.64~\mathrm{kHz}$, respectively.
    (d) Gaussian-fit FWHM in units of $\Gamma_{T_2}$; colors match (c).
    The red curve is the constant-drive Bloch power-broadening reference;
    the dotted line marks unity, with
    $\Gamma_{T_2}=\EffectiveCoherenceFwhm~\mathrm{kHz}$,
    using the primary qubit's coherence calibration.
    The $\beta=-0.22$ acquisition combines the derivative quadrature
    with a smoothed midpoint inversion, so this comparison does not
    isolate the DRAG contribution or establish an optimum in $\beta$.}
    \label{fig:ac-stark-correction-maps}
    \label{fig:main-ac-stark-shifts}
\end{figure}

For comparison, the weak-drive three-level AC-Stark shift of a constant drive
is $\delta f_{01}\simeq-\left(\Omega_0/2\pi\right)^2/
[2(\alpha/2\pi)]$~\cite{PhysRevA.76.042319,schneider_multilevel_2018}.
With $\alpha/(2\pi)=-237~\mathrm{MHz}$ for this qubit, the shift reaches
$5.3~\mathrm{MHz}$ at $\Omega_0/(2\pi)=50~\mathrm{MHz}$.  For the
$20~\mu\mathrm{s}$ root-Lorentzian envelope with $c=0.001$, the phase-average
shift is smaller by $\langle f^2\rangle_t=0.001570$ and reaches only
$8.2~\mathrm{kHz}$; the echo sign change does not cancel this quadratic term.
Fig.~\ref{fig:ac-stark-correction-maps} instead tests a derivative-quadrature
waveform on the second qubit. The acquisition
uses $I+iQ$ with $Q=-\beta\dot I/(2\pi\alpha)$, where
$\alpha=237~\mathrm{MHz}$ is the signed value supplied to the
waveform generator and time is in microseconds. Both the derivative
quadrature and the midpoint smoothing change between these sequential runs;
no accumulated-phase correction is applied. The depletion feature persists
across the displayed $10$--$60~\mathrm{MHz}$ range in both maps.
The RMS center motion about each run's mean is $6.78$ and
$2.64~\mathrm{kHz}$ for $\beta=0$ and $-0.22$, respectively. This lower
scatter characterizes the combined smoothed-midpoint and DRAG waveform;
the comparison does not isolate the derivative-quadrature contribution.
The absence of resolved three-state readout precludes a measured
leakage-suppression claim.

Fig.~\ref{fig:main-ac-stark-shifts}(c) resolves the measured center
variations directly on a kilohertz scale. Panel (d) compares FWHM values
extracted from the same Gaussian dip estimator with the steady-state Bloch
reference $\Gamma=\Gamma_{T_2}\sqrt{1+\Omega^2T_1T_2}$, using the primary
spectroscopy values $T_1=\MeasuredTOne~\mu\mathrm{s}$ and
$T_2=\EffectiveTTwo~\mu\mathrm{s}$. This is a primary-qubit coherence
reference, rather than a separate calibration of the second qubit.
Both measured linewidths
increase with drive;
the smaller center fluctuations for $\beta=-0.22$ therefore do not imply
a uniformly narrower resonance. Numerical accumulated-phase compensation
is discussed separately in the Supplemental Material~\cite{SupplementalMaterial}.
